\chapter{Literature Review}
\section{Robot Perception}
\section{Robot Localization}
\section{Monocular Depth Estimation}
Early approaches to MDE relied on hand-crafted geometric cues and probabilistic graphical models,
such as Markov random fields, to infer depth from superpixel features
\todo{cite classical MDE work, e.g. Saxena et al.}. The advent of deep learning shifted the field
toward end-to-end CNN regression, where an encoder-decoder network is trained to directly
predict a dense depth map from an RGB input, supervised by ground-truth depth from LiDAR
or stereo datasets \todo{cite Eigen et al. and follow-up encoder-decoder work}. More recently,
vision transformer (ViT) backbones have been shown to outperform convolutional architectures
on this task, since the global self-attention mechanism of a transformer captures long-range
spatial context that is useful for resolving scale ambiguity, whereas CNNs are inherently
biased toward local receptive fields \todo{cite DPT / MiDaS or similar ViT-based depth work}.

A second axis along which this literature divides is whether a model is trained to generalize
across scenes at all. Models fine-tuned on a single dataset with fixed camera intrinsics, such
as NYU Depth v2 or KITTI, can be trained to regress depth in meters, but transfer poorly beyond
the domain they were fitted to. Zero-shot models instead train on a mixture of heterogeneous
sources and trade metric output for breadth of generalization
\todo{cite MiDaS, Depth Anything}. Chapter~\ref{ch:mde} takes up the consequences of that
trade-off for the Cyclops rover.

%%% Local Variables: ***
%%% mode: latex ***
%%% TeX-master: "thesis.tex" ***
%%% End: ***
